\documentclass[11pt]{article}
% NSF PAPPG Ch. 2 SS B.2: font >=11pt, margins >=1in, <=6 lines/inch.
% No hyperlinks in Project Description body (Research.gov constraint).

\usepackage{bm}
\usepackage{color}
\usepackage{hyperref}
\usepackage{graphicx}
\usepackage{wrapfig}
\usepackage{subcaption}
\usepackage{vmargin}
\usepackage{mathrsfs}
\usepackage{float}
\usepackage{tabularx}
\usepackage{booktabs}
\usepackage{longtable}
\usepackage{ltablex}
\usepackage[svgnames]{xcolor}
\usepackage[most]{tcolorbox}
\usepackage{enumitem}
\usepackage{titlesec}
\usepackage{caption}
\usepackage{colortbl}
\usepackage[numbers,square,sort&compress]{natbib}
\usepackage{pgfgantt}
\usepackage{tikz}
\linespread{0.935}

\definecolor{mycmykcolor}{cmyk}{0.45,0.30,0.20,0}
\definecolor{secondary}{cmyk}{0,0,0,0.20}
\definecolor{roonecol}{HTML}{8C9EB2}
\definecolor{rootwocol}{HTML}{B08A72}
\definecolor{roothreecol}{HTML}{7A9985}
\definecolor{lightgray}{HTML}{E8E8E8}
\definecolor{figuregreen}{HTML}{A8D5C2}
\definecolor{figureblue}{HTML}{A8C8E8}
\definecolor{figuregray}{HTML}{D9D9D9}
\definecolor{figureborder}{HTML}{156082}

\pagestyle{empty}
\titleformat{\section}[runin]{\normalfont\bfseries}{\thesection}{1em}{}[:]
\titlespacing*{\section}{0pt}{\baselineskip}{1ex}
\titleformat{\subsection}[runin]{\normalfont\bfseries}{\thesubsection}{1em}{}[:]
\titlespacing*{\subsection}{0pt}{\baselineskip}{1ex}

\bibliographystyle{unsrtnat}
\hypersetup{hidelinks,colorlinks=false,urlcolor=black}
\hypersetup{draft}
\setpapersize{USletter}
\setmarginsrb{1in}{1in}{1in}{1in}{0pt}{0mm}{0pt}{0mm}
\pretolerance=5000
\tolerance=9000
\graphicspath{{images/}}

\newcommand{\pseudodot}{{\lower 2.4pt\hbox{$\cdot$}}}
\newcommand{\colorsquare}[1]{\tikz[baseline=-0.6ex]\draw[fill=#1, draw=none] (0,0) rectangle (1ex,1ex);}
\newcommand{\framesquare}[1]{\tikz[baseline=-0.6ex]\draw[fill=none, draw=#1] (0,0) rectangle (1ex,1ex);}


\newcommand{\customsection}[1]{%
  \titleformat{\section}[block]{\normalfont\bfseries}{\thesection}{1em}{}[]%
  \section{#1}%
  \titleformat{\section}[runin]{\normalfont\bfseries}{\thesection}{1em}{}[:]%
}

\begin{document}

% -----------------------------------------------------------------------
% GOALS TCOLORBOX

% current title: Social and Physical Flood Criticality in Heritage Districts: A Coupled Social-Hydrodynamic Framework

% potential title: Unraveling the Preservation-Urgency Paradox: Flood Risk and Collective Action in URM Heritage District
% -----------------------------------------------------------------------

% \begin{tcolorbox}[colback=lightgray, colframe=black, boxrule=1pt, coltext=black]
% \textbf{Project goals:}
% Heritage districts, dense with older unreinforced masonry (URM) are socio-technical infrastructure systems whose flood resilience
% depends on the state of their physical assets, the organizational capacity to act on them,
% and the service infrastructure that connects them.
% We couple social science, hydrodynamic modeling, and spatial data science to
% investigate how the preservation-urgency paradox operates at individual and district
% scales in URM heritage districts, test whether the spatial configuration of hardening
% decisions has nonlinear consequences for collective flood outcomes, and determine
% whether collective action exclusion and physical flood criticality co-locate within
% districts, translating findings into actionable toolkit resources.
% \textbf{RQ1} (Napolitano, Busse, Chi) asks how the preservation-urgency paradox
% operates at both the individual building and district scales, and which buildings and owners are drawn into collective action versus acting independently.
% \textbf{RQ2} (Busse, Napolitano) asks whether the spatial configuration of
% pre-disaster hardening decisions across a heritage district produces nonlinear
% effects on district-scale flood vulnerability, and whether individual building states
% produce disproportionate collective consequences at identifiable thresholds.
% \textbf{RQ3} (Chi, Napolitano, Busse) asks whether collective action exclusion in
% heritage district flood recovery has a spatial signature corresponding to where
% individual hardening decisions produce disproportionate collective flood consequences.
% The defining intellectual contribution is the first empirical test of whether
% social exclusion from collective action and physical flood criticality are coupled
% phenomena in heritage districts, producing findings that give communities a
% spatially targeted basis for pre-disaster collective action investment (Fig.~\ref{fig:framework}).
% \end{tcolorbox}

% -----------------------------------------------------------------------
% OPENING WRAPFIGURE
% -----------------------------------------------------------------------
\begin{wrapfigure}{r}{0.55\textwidth}
\includegraphics[width=0.55\textwidth]{images/ISap_fig1_may26.png}
\captionsetup{justification=centering,singlelinecheck=false,font=small}
\vspace{-0.9cm}
\caption{Coupled social-hydrodynamic framework linking community decision-making under the preservation urgency paradox to district-scale flood vulnerability outcomes in URM heritage districts. Intervention typologies (e.g., \protect\colorsquare{figuregreen} wet floodproofing, \protect\colorsquare{figureblue} dry floodproofing, \protect\colorsquare{figuregray} no action) and collective action participation statuses (e.g., \protect\framesquare{figureborder} coordinated, \protect\framesquare{figuregray} excluded) are defined empirically in RO1.}
\vspace{-0.2cm}
\label{fig:framework}
\end{wrapfigure}

% -----------------------------------------------------------------------
\section{Status quo}
% -----------------------------------------------------------------------

The United States faces intensifying flood hazards that disproportionately threaten
heritage-rich inland communities whose unreinforced masonry (URM) building stock was
constructed before modern flood resilience standards existed~\cite{Xu2024,QIANG2019,probfut2026}.
Predicted annual flood damage ranges from \$180 to \$500 billion across 94\% of
continental states~\cite{JEC2024FloodingCosts}, with inland non-metropolitan community
centers experiencing compounded vulnerabilities such as geographic isolation, limited
institutional infrastructure, constrained financial resources, and building stock whose
historic significance imposes preservation obligations that standard disaster assistance
frameworks are not designed to support~\cite{WattsTompkins2022,ClimateCentral2024,wilson2021flood,zhorov2024floodprints}.
When floods strike these communities, the irreplaceable cultural assets that define
local identity and anchor tourism-dependent economies face permanent loss through both
direct damage and cascading business interruption that persists even when physical
damage is limited~\cite{OlderSmallerBetter2014,untapped2017,firststreet2026}.
The scale of this threat is illustrated by Montpelier, Vermont, where 80\% of the
National Register Historic District sustained flood damage in July
2023~\cite{NWSPaducah2023,dupigny2002climate}, and Marshall, North Carolina, where
Hurricane Helene devastated a rural Appalachian heritage district in September
2024~\cite{npr_hurricane_helene_2024}.

The fundamental difficulty is systemic incompatibility at two coupled scales.
Heritage districts are socio-technical infrastructure systems: physical assets operated
through human organizational agency where functional recovery requires both built
infrastructure and organizational capacity~\cite{NIST2024SP1310}.
At the building scale, structures are hardened through National Park Service (NPS)-guided
pre-disaster intervention to reduce vulnerability while maintaining historic integrity,
demanding time, resources, and institutional coordination these communities chronically
lack~\cite{NPS2021,SeronSuhoothi2017,Jigyasu2024,joshi2021community}.
At the community scale, the district's combined built fabric, organizational capacity,
and service infrastructure determines flood inundation, velocity, and duration, yet
the relationship between individual preservation decisions and system-level outcomes
has never been characterized~\cite{NIST2024SP1310,vojinovic2016holistic}.
The NPS flood adaptation guidelines identify intervention options at both scales,
from individual building hardening such as wet and dry floodproofing to district-scale
physical modifications such as riparian vegetation management and drainage
improvements~\cite{NPS2021}, yet no existing work characterizes how these approaches
compare in their collective flood vulnerability outcomes, or what organizational and
network structures determine which approach a community pursues.
The ASCE-JSCE infrastructure resilience framework formalizes this requirement,
treating organizational capacity as a co-equal determinant of service restoration
alongside physical performance~\cite{davis2023}.
The two scales are coupled through the preservation-urgency paradox: tension between
rapid disaster imperatives and technically-informed heritage preservation that causes
communities to diverge from NPS-recommended infrastructure states, with unquantified
consequences for district-scale flood
vulnerability~\cite{maalouf2025a,maalouf2025b} (Fig.~\ref{fig:paradox}).
Existing research has advanced technical understanding of flood impacts on heritage
districts and community recovery approaches, but treats physical vulnerability and
social decision-making as separate
domains~\cite{sesana2021climate,NPS2021,Arrighi2024,castro2021community,gibson2019engagement,briz2022prioritization,fatoric2023towards,kim2021implementation}.
NIST SP 1310's recommended future research agenda explicitly identifies social
organizational and multidisciplinary social-technical computational models as unmet
needs, and calls for extending the framework to hazards beyond earthquakes alongside
the FEMA/BSSC Functional Recovery Planning
Committee~\cite{NIST2024SP1310,NIBSFunctionalRecovery2023}; this research pursues
that integration for flood hazards and heritage districts.

\begin{wrapfigure}{r}{0.7\textwidth}
\vspace{-0.65cm}
\includegraphics[width=0.7\textwidth]{images/ISap_fig2_may26.png}
\captionsetup{justification=centering,singlelinecheck=false,font=small}
\vspace{-0.8cm}
\caption{The preservation-urgency paradox operates at both pre-disaster
(preparedness failure) and post-disaster (recovery divergence) temporal scales,
and at both individual building and district spatial scales, producing measurable
divergence from NPS-recommended infrastructure states whose consequences for
community-scale flood vulnerability this research quantifies. Intervention typologies
(e.g., \protect\colorsquare{figuregreen} wet floodproofing,
\protect\colorsquare{figureblue} dry floodproofing,
\protect\colorsquare{figuregray} no action) and collective action participation
statuses (e.g., \protect\framesquare{figureborder} coordinated,
\protect\framesquare{figuregray} excluded) are defined empirically in RO1.}
\vspace{-0.45cm}
\label{fig:paradox}
\end{wrapfigure}

While recent work couples social vulnerability with building-scale flood recovery
trajectories for residential buildings~\cite{mazumder2025modeling}, no study has
coupled community-scale hydrodynamic modeling with empirically documented heritage
district intervention decision patterns to characterize how individual building hardening,
district-scale physical interventions, and mixed approaches compare in their effects
on district-scale flood vulnerability, or what organizational and network structures
explain why communities default to fragmented individual hardening rather than
coordinated district-scale intervention where technical expertise identified it as
an option.
This research addresses the ISP program's core interest in how infrastructure systems
function under extreme events and how human decision-making shapes system resilience,
by producing the first empirical comparison of building-level and district-scale
intervention approaches in URM heritage districts and the first characterization of
the organizational dynamics that produce coordination failure around district-scale
intervention, directly advancing the socio-technical systems framing established in
NIST SP 1310~\cite{NIST2024SP1310}.
% -----------------------------------------------------------------------
% STATUS QUO WRAPFIGURE
% -----------------------------------------------------------------------

\vspace{-0.4cm}
% -----------------------------------------------------------------------
\subsection{Intellectual merit}
% -----------------------------------------------------------------------
This research advances infrastructure systems science by establishing how the
preservation-urgency paradox operates at individual and district scales, testing
how individual building hardening, district-scale physical intervention, and mixed
approaches compare in their effects on collective flood vulnerability, and determining
what organizational and network structures produce the coordination failures that
prevent communities from pursuing district-scale intervention.
Building on preliminary identification of this paradox~\cite{maalouf2025a,maalouf2025b},
we challenge conventional assumptions about community capitals in heritage preservation
recovery~\cite{MurzynKupisz2013,HistoricEngland2023}.
While self-organized, community-driven approaches have been central to disaster recovery
by drawing on community capitals to assess needs, allocate resources, and coordinate
recovery~\cite{imperiale2016experiencing}, our research investigates how collective action
participation is structured, with preliminary findings suggesting network position rather
than heritage significance as the primary determinant~\cite{maalouf2025a,maalouf2025b},
and whether the network structures that exclude properties from collective action also
explain why communities default to fragmented individual building hardening rather than
coordinated district-scale intervention such as riparian vegetation management or drainage
improvements~\cite{Hou2023,UNESCO2024}.
This investigation will contribute to fundamental social and engineering sciences by
advancing understanding of how collective action structure shapes infrastructure system
resilience in flood-affected heritage districts~\cite{moya2022approach,aldrich2015social,cleaver2002reinventing}.
The research questions that organize this work are:
\newpage
\begin{itemize}[noitemsep]
  \item \textbf{RQ1 (Napolitano, Busse, Chi):} How does the preservation-urgency paradox 
  operate at both the individual building and district scales in URM heritage districts, 
  and which buildings and owners are drawn into collective action around building-level 
  hardening or district-scale intervention versus acting independently?
  \item \textbf{RQ2 (Busse, Napolitano):} How do individual building hardening, 
  district-scale physical interventions, and mixed approaches compare in their effects 
  on district-scale flood vulnerability in URM heritage districts, and does any 
  intervention approach produce threshold behavior in collective flood outcomes?
  \item \textbf{RQ3 (Chi, Napolitano, Busse):} Among buildings and owners excluded 
  from collective action in heritage district flood recovery, what organizational and 
  network structures explain that exclusion, and do those same structures explain why 
  district-scale physical intervention was not pursued where technical expertise 
  identified it as an option?
\end{itemize}
\vspace{-0.3cm}
The three RQs are coupled in a feed-forward structure that constitutes the primary
intellectual contribution: the collective action participation patterns, intervention
typology classifications, and awareness of district-scale options documented in RO1
provide the spatially explicit inputs for the hydrodynamic modeling in RO2, and the
RO2 outputs comparing intervention approaches are combined with the collective action
exclusion patterns and network structures from RO1 to test whether the organizational
dynamics that produce social exclusion also explain the failure to pursue district-scale
intervention, the question at the core of RQ3.
\vspace{-0.3cm}
\begin{figure}[H]
\centering
\includegraphics[width=0.9\textwidth]{table.png}
\vspace{-0.3cm}
\caption{Summary of research gaps, objectives, approaches, and outcomes.\label{fig:research_summary}}
\end{figure}

% -----------------------------------------------------------------------
\section{Current state-of-the-art}
% -----------------------------------------------------------------------

Research on disaster decision-making in flood-vulnerable heritage communities spans three domains
including technical preservation of URM heritage structures under flood hazards,
community-based disaster recovery in heritage-rich environments, and infrastructure
systems modeling of flood vulnerability.
Each domain has advanced independently, and no existing work integrates all three.
The integration gap, specifically the absence of a coupled social-hydrodynamic framework
linking community decision-making patterns to district-scale flood vulnerability in
these communities, is the scientific opportunity this research addresses.

URM structures present distinct technical challenges under flood loading;
retrofitting requires balancing structural intervention with preservation of historic
fabric, and standard damage models substantially underestimate failure risk because
prolonged inundation compromises lime mortar and high-porosity brick in ways not
accounted for in flood exposure assessments, with wall collapse fragility nonlinear
above approximately 90~cm inundation
depth~\cite{SeronSuhoothi2017,xiao2021flood,castillo2026moisture,korswagen2025experimental,moaiyedfar2025collapse}.
The NPS flood adaptation guidelines provide the federal typology for heritage
intervention at both building and district scales~\cite{NPS2021}, though empirical
evidence from inland flood recovery contexts suggests that compliance requirements
often exceed owner capacity during crisis conditions, leading many property owners
to forgo preservation programs despite potential
benefits~\cite{maalouf2025b,GAO2025,Poling2024}.
Post-flood damage extends beyond structural losses to include diminished cultural
significance that standard assessments do not
capture~\cite{Arrighi2024,sesana2021climate}.
Flood risk frameworks exist at global and individual-asset scales but have not been
applied to inland U.S. heritage districts where collective action patterns determine
which buildings receive intervention~\cite{carmichael2023riverine,figueiredo2024flovri,vojinovic2016holistic}.
Existing work examines individual building performance without characterizing how
building-level decisions produce district-scale outcomes, the gap this proposal
addresses through NIST SP 1310 socio-technical framing~\cite{NIST2024SP1310}.

District-scale hydrodynamic modeling estimates flood extent, depth, and velocity
from terrain, drainage, and building morphology interactions, with 1D/2D approaches
well established for urban
environments~\cite{usgs_effects,sanders2015hydrodynamic,fan2017coupled,leandro2009comparison,mustafa2023gis}.
Nature-based solutions such as riparian vegetation management and drainage
improvements are increasingly modeled within HEC-RAS 2D by adjusting Manning's
$n$ to represent increased flow resistance from vegetation and modifying terrain
to represent drainage and regrading interventions, with established parameter
ranges for afforestation, wetland creation, and floodplain
reconnection~\cite{wang2019evaluating,modellingnbs2026}; no existing work applies
this approach to compare building-level and district-scale intervention approaches
in heritage districts.
Scenario-based assessments treat building stock as static input without the social
decision-making that determines which configurations are
realized~\cite{zhu2019modeling,bsi128,sanders2015hydrodynamic}.
Agent-based model (ABM)-coupled frameworks focus on evacuation and post-event
response rather than pre-disaster collective action; no existing socio-hydrodynamic
framework models this decision class~\cite{aerts2020abm,anshuka2022abm}.

Community-based disaster research establishes that social capital (bridging connections and linking relationships with institutions)
shape vulnerability and recovery trajectories in heritage-rich environments by enabling resource
mobilization, coordination, and knowledge
exchange~\cite{aldrich2015social,MurzynKupisz2013,imperiale2016experiencing}.
In heritage specifically, community organization and participatory methods
determine who is involved in decisions, which owners receive support, and which
preservation approaches are
adopted~\cite{castro2021community,gibson2019engagement,joshi2021community,Jigyasu2024}.
Preliminary work at Montpelier and Marshall finds that preservation success correlates
with network access rather than heritage significance or recovery need, a finding
that challenges the Putnam-Aldrich framework in which higher social capital
produces better recovery outcomes regardless of how that capital is
distributed~\cite{aldrich2015social,putnam2015bowling}: if network position
rather than need determines who receives support, social capital
functions as a sorting mechanism that concentrates recovery resources among
already-connected owners rather than as a collective good.
Whether the same network structures that produce exclusion also explain why
communities default to fragmented individual building hardening rather than
coordinated district-scale intervention, even where technical expertise identified
it as an option, remains untested (Fig.~\ref{fig:research_summary}).
\vspace{-0.3cm}
% -----------------------------------------------------------------------
\section{Pilot sites and preliminary work}
% -----------------------------------------------------------------------

This research is grounded in two sites selected for their analytic contrast and
existing PI research relationships.
Both sites are NPS-designated heritage districts that experienced recent catastrophic
flooding and have established data access through prior work; they differ meaningfully
on the dimensions most relevant to the research questions: hydraulic regime (valley
river versus confined mountain channel), institutional capacity (state capital with
preservation infrastructure versus small rural town with minimal formal support), and
federal program access.
Findings from each site are reported as site-specific results; convergence across
sites is treated as suggestive evidence warranting future comparative work at larger
scale, not as proof of generalizability.
Information on the communities is provided in Tab.~\ref{tab:community_stats}.
\vspace{-0.3cm}
\begin{table}[h]

\centering
\footnotesize
\caption{Statistics of the two study communities~\cite{uscensus_acs5year}}
\label{tab:community_stats}
\vspace{-0.8cm}
\begin{tabularx}{\textwidth}{@{}l*{9}{>{\centering\arraybackslash}X}@{}}
\toprule
\textbf{Community} & \textbf{Pop.} & \textbf{\% White} & \textbf{Med. Age} & \textbf{Med. Inc.} & \textbf{\% Poverty} & \textbf{Med. Prop.} & \textbf{Prop. Tax} & \textbf{\% Home} & \textbf{\% Mort.} \\
\midrule
Montpelier, VT & 8,030 & 90 & 44 & \$76,143 & 6.4 & \$326,800 & \$2,154 & 56 & 59.7 \\
Marshall, NC   & 616   & 87 & 42 & \$49,700  & 22.7 & \$228,400 & \$183  & 60 & 65 \\
Nationwide     & --    & 89 & 39 & \$75,149  & 12.5 & \$281,900 & \$2,111 & 65 & 65 \\
\bottomrule
\end{tabularx}
\footnotesize{Note: Pop.\ = Population; Med.\ = Median; Inc.\ = Income; Prop.\ = Property value;
\% Home = Homeownership rate; \% Mort.\ = Percentage of homes with mortgage}
\end{table}
\normalsize
\vspace{-0.3cm}
\begin{wrapfigure}{r}{0.4\textwidth}
\vspace{-0.4cm}
\includegraphics[width=0.4\textwidth]{images/vt.pdf}
\captionsetup{justification=centering,singlelinecheck=false,font=small}
\vspace{-0.9cm}
\caption{Montpelier 2023 flooding.}
\label{fig:vt}
\vspace{-0.2cm}
\end{wrapfigure}
Montpelier's baseline is enriched by a structural dataset acquired through an
NHERI RAPID Graduate Scholar grant and archived on DesignSafe~\cite{maalouf2024data},
providing material composition, roof form, wall thickness, foundation type,
upper-story damage state, and wall geometry where Street View cannot resolve
sufficient detail~\cite{maalouf2025a}.
Feature importance analysis of 271 heritage structures within this dataset revealed
that commercial use, high property values, and larger building areas were the strongest
predictors of faster recovery permit initiation following disaster, while residential
and religious or cultural properties showed significantly longer delays~\cite{maalouf2025a},
establishing the quantitative equity baseline against which RQ3's network and
organizational analysis will be evaluated (Fig.~\ref{fig:vt}).
\begin{wrapfigure}{l}{0.4\textwidth}
\vspace{-0.4cm}
\includegraphics[width=0.4\textwidth]{images/nc.pdf}
\captionsetup{justification=centering,singlelinecheck=false,font=small}
\vspace{-0.6cm}
\caption{Marshall, NC: Hurricane Helene flooding aftermath, September 2024.}
\label{fig:nc}
\vspace{-0.2cm}
\end{wrapfigure}
For Marshall, a December 2024 field mission by PI Napolitano's group produced
drone photogrammetry covering approximately 30\% of the district along with
building condition assessments and preliminary permit records~\cite{maalouf2026data};
Subtask 1.1 will expand drone coverage to the full district.
Preliminary analysis identified four intervention windows governing heritage-conscious
recovery outcomes: material salvage (days 0--10), expert consultation (days 7--20),
funding mobilization (days 10--25), and implementation commitment (days 20--33),
with timeline flexibility the strongest predictor of NPS-consistent
outcomes~\cite{maalouf2025b} (Fig.~\ref{fig:nc}).

% maggie: please review the paragraph below for accuracy and add citations for
% your HEC-RAS 2D experience and any published work at comparable sites.
% %MAGGIE please add citations

Co-PI Busse brings established district-scale hydrodynamic modeling capability using
HEC-RAS 2D, with experience spanning riverine fluvial flooding and urban surface
inundation at the neighborhood scale, the modeling context this proposal requires.
The two pilot sites present hydraulically distinct systems: Montpelier's mid-size
Vermont valley river (moderate gradients, broad floodplain, dense urban fabric) versus
Marshall's confined Appalachian mountain channel (rapid stage rise, high velocities,
limited buffering), allowing threshold behavior to be tested across mechanistically
different flood regimes and strengthening transferability of findings.

% guangqing: please review the paragraph below for accuracy and add any prior
% work from your lab on collective action, social networks, or rural community
% resilience that is relevant here.

The social science foundation for RO1 builds on two preliminary studies and
Co-PI Chi's prior work on collective action in rural governance~\cite{li2023collective}.
In Montpelier, thematic coding \cite{boyatzis1998transforming} of 15 semi-structured stakeholder interviews in
ATLAS.ti, combined with feature importance analysis of recovery trajectories across
271 heritage structures, revealed that building characteristics and socioeconomic
factors produced divergent recovery trajectories, and that institutional arrangements
privileged certain recovery pathways over others regardless of heritage significance
or recovery need~\cite{maalouf2025a}.
In Marshall, a mixed-methods analysis of 21 heritage buildings and interviews with
16 property owners conducted 10 weeks post-disaster showed that in the absence of
formal preservation support, social networks functioned as improvised preservation
infrastructure, but that preservation success correlated with network access rather
than heritage significance or recovery need~\cite{maalouf2025b}.
The validated mixed-methods framework from these studies will be applied at both
sites in RO1, with Chi leading social network and institutional analysis that
extends the thematic findings into a spatially explicit exclusion map.
\vspace{-0.4cm}
% -----------------------------------------------------------------------
\section{Research plan}
% -----------------------------------------------------------------------
The research plan is organized around three research objectives corresponding to the
three research questions.
RO1 (Napolitano, Busse, Chi; 4 subtasks) establishes the coupled physical and social
data foundation at both sites using a shared NPS-grounded intervention typology spanning
both building-level and district-scale interventions.
RO2 (Busse, Napolitano; 3 subtasks) tests how individual building hardening,
district-scale physical interventions, and mixed approaches compare in their effects
on district-scale flood vulnerability, and whether any intervention approach produces
threshold behavior in collective flood outcomes.
RO3 (Chi, Napolitano, Busse; 3 subtasks) characterizes the organizational and network
structures that produce collective action exclusion and explains why district-scale
physical intervention was not pursued where technical expertise identified it as an
option, and translates findings into a field-tested, federally-informed toolkit
co-developed with the Association of Preservation Technology (APT DRI) which PI
Napolitano leads and the Department of the Interior.

\begin{tcolorbox}[colback=roonecol!60,colframe=black,boxrule=1pt,coltext=black,
  left=1mm,right=1mm,top=1mm,bottom=1mm]
\textbf{Research Objective 1 (RO1, leads: Napolitano, Busse \& Chi): Characterize
how the preservation-urgency paradox operates at both the individual building and
district scales, documenting collective action participation and classifying
infrastructure states at both building and district scales to produce the coupled
dataset that RO2 and RO3 require.}
\end{tcolorbox}

\textbf{Subtask 1.1: NPS intervention typology and collective action classification (lead: Napolitano).}
Using the NPS flood adaptation guidelines as a classification framework, each structure
in Montpelier and Marshall will be assigned an intervention type reflecting its
pre-disaster condition relative to NPS recommendations, including wet floodproofing,
dry floodproofing, elevation, no action, or non-compliant modification at the building
scale, and riparian vegetation management, drainage improvements, and regrading at the
district scale.
Classification draws on Street View imagery, which provides complete building-level
coverage at both sites, as the primary source for visual inspection and typology
assignment; building footprints are drawn from OpenStreetMap.
For buildings where Street View cannot resolve key structural details, laser scan data
(Montpelier) and drone photogrammetry (Marshall) provide enrichment, adding roof
condition, material specifics, upper-story damage state, and wall geometry.
Building records and permit data supplement imagery review for intervention history.
District-scale intervention states are documented through organizational records,
permit data, and field observation; Montpelier's post-2023 riparian vegetation work
provides one empirically observed district-scale configuration, while Marshall's
district-scale intervention state reflects no observed action.
Each building's record will also assign a collective action participation status
using a two-step construction that distinguishes exclusion from voluntary
non-participation.
First, a reach indicator is derived from organizational records: emails, meeting
attendance lists, event photographs, and program documentation establish which
owners were contacted by any collective action effort.
Second, a response indicator is derived from interviews: of those reached, whether
the owner participated or declined.
This produces three analytically distinct categories: coordinated (reached and
participated), exit (reached but declined), and excluded (no evidence of contact
in organizational records).
RQ3 uses the excluded category specifically; the exit category is retained as a
comparison group for the social network analysis~\cite{scott2012social} in Subtask 3.1.
A shared building-level identifier links each record across the imagery review,
permit records, Chi's interview dataset, and Busse's hydrodynamic model inputs,
enabling the RO3 network and organizational analysis connecting collective action
exclusion patterns with intervention approach outcomes from RO2.
The typology uses one-hot encoding across observable treatment categories, capturing
multi-treatment buildings without forcing a single classification; categories
requiring structural access are resolved through permit records and interviews.
Jennifer Wellock (co-author of the NPS flood adaptation guidelines~\cite{NPS2021},
previous collaborator of PI Napolitano, see Letters of Collaboration)
will review the completed typology for consistency with current federal practice.

\textbf{Subtask 1.2: District-scale hydrodynamic site characterization (lead: Busse).}
Busse will develop site-specific hydrodynamic baselines for Montpelier and Marshall
by using high resolution digital elevation models with OpenStreetMap (OSM) vector
features (building footprints, channel alignment) to define model geometry.
Baselines will establish the channel geometry, building footprints, elevation, and
permeability conditions that Subtask 2.1 will use as the foundation for both
building-level and district-scale infrastructure configuration development; where
available, laser scan and drone photogrammetry data enrich building geometry inputs
with structural depth beyond what OSM provides.
The hydraulic contrast between the two sites presents a genuine parameterization
challenge: Montpelier's mid-size Vermont valley river system and Marshall's steep
confined Appalachian mountain channel require distinct approaches to boundary
conditions, flood event type, and calibration data, which Busse will address
through site-specific model development before configurations are introduced in RO2.

% maggie: please review the paragraph below for accuracy and add citations for
% HEC-RAS 2D and any USGS lidar data products. %MAGGIE please add citations

Site-specific HEC-RAS 2D models~\cite{HECRAS2D} will be developed using USGS lidar
terrain data (QL1 at 0.35~m for Montpelier, QL2 at 0.7~m for Marshall)~\cite{USGS3DEP}.
15-minute stage and discharge records from USGS gauges at the Winooski
river~\cite{USGS_Winooski_Waterbury,USGS_Winooski_Montpelier,USGS_Winooski_EssexJunction}
will provide boundary conditions for the models.
Models will be calibrated against the July 2023 and 2024 floods at Montpelier and
Hurricane Helene (September 2024) at Marshall, using gauge records, peak stage
observations, and satellite-derived inundation data, with validation against
high water marks~\cite{Olson2025_VT2023,Olson2025_VT2023_SIR,USGS_FEV_VT2024,USGS_FEV_Helene2024},
FEMA flood maps~\cite{FEMA_FIRM_Marshall,FEMA_FIRM_Montpelier,FEMA_DR4720_VT,FEMA_DR4782_VT,FEMA_DR4827_NC},
and post-event survey data.
The two flood regimes require distinct parameterization: Montpelier's rainfall-driven
riverine system and Marshall's hurricane-driven confined mountain channel differ
in stage rise rate, velocity, and floodplain interaction.
Busse will develop site-specific boundary conditions and roughness schemes for each.
For Marshall's steep confined channel, where flow is sensitive to sub-meter
terrain variation, calibration will confirm that 1~m resolution is adequate for
characterizing building-level flood depth differences; if calibration residuals
indicate resolution-driven error, the model domain will be refined at critical
locations using the QL2 0.7~m terrain data.
Baseline outputs are 1~m resolution raster grids of inundation depth, flow velocity,
and water surface elevation; each building in the OSM footprint is resolved at this
grid, enabling Napolitano to assign NPS typology hardening states as building-level
model parameters (dry floodproofing as blocked obstructions, wet floodproofing and
no-action as flow-through with adjusted Manning's~$n$) and district-scale intervention
states as terrain edits and Manning's~$n$ adjustments for riparian and drainage
modifications for RO2 configuration runs.

\begin{wrapfigure}{r}{0.36\textwidth}
\vspace{-0.4cm}
\includegraphics[width=0.36\textwidth]{images/josh.jpg}
\captionsetup{justification=centering,singlelinecheck=false,font=small}
\vspace{-0.4cm}
\caption{Preliminary community conversations in Marshall, NC (IRB-approved).}
\label{fig:josh}
\vspace{-0.4cm}
\end{wrapfigure}
\textbf{Subtask 1.3: Social science data collection (Chi).}

Chi will lead mixed-methods data collection on community decision-making shaping
disaster response in heritage districts at both sites (Fig.~\ref{fig:josh}).
Archival records identify which buildings participated in collective action; interviews
explain mechanisms of participation and exclusion; and quantitative data on building
characteristics, socioeconomic factors, and community capital~\cite{uscensus_acs5year,li2023collective}
provides the feature space for Subtask 1.4.
Interviews use stratified purposive and snowball sampling~\cite{maxwell1996methods}
with local authorities, community leaders, heritage experts, property owners, and residents.
Protocols focus on decision points within the four intervention
windows~\cite{maalouf2025a,maalouf2025b}, factors enabling NPS-consistent preparation,
collective action participation including why owners were or were not engaged and
what network structures shaped their access, awareness of district-scale intervention
options including whether technical experts communicated those options and through
what channels, and what structural or organizational barriers prevented coordinated
district-scale action where it was not pursued.
At Marshall, organizational records from the Downtown Marshall Association, including
emails and meeting notes from recovery coordination efforts and event photographs from
post-Helene community gatherings, establish reach independent of interview recall;
interviews validate and add nuance to the record-derived classifications rather than
serving as the primary source.
IRB protocol submission in month 1 will confirm that organizational record collection
is within scope; PI Napolitano will expand coverage beyond the existing exemption
determination (STUDY00024770) if the organizational records protocol requires it.
All interview and archival records will be linked to the shared building identifier
established in Subtask 1.1 to enable the network and organizational analysis in
Subtask 3.1.
We anticipate approximately 50--60 interviews in Montpelier and 40--50 in Marshall,
with sampling continuing until data saturation is reached~\cite{maxwell1996methods}.
IRB approval will be obtained prior to data collection; PI Napolitano has existing
IRB exemption determination for prior work at both sites.

% guangqing: please review the paragraph below, add your preferred citations
% for purposive/snowball sampling and thematic analysis methods, and confirm
% the target interview counts per site.
% %GUANGQING please add: target sample size per site and any lab citations

Protocols build on the validated instruments from the Montpelier and Marshall
preliminary work~\cite{maalouf2025a,maalouf2025b}, extended to capture whether
building owners were reached by district-level collective processes: whether
organizations contacted them, what was offered, whether they participated, what
structural barriers prevented engagement where they did not, and whether
district-scale intervention options were communicated to property owners and
community leaders or remained within technical and engineering networks.
Interviews will be audio-recorded, transcribed, and coded in NVivo with codes
for collective action exposure, engagement response, structural barrier type, and
awareness of district-scale intervention options.
Building permit data for Montpelier (2023 events) are already secured; records
requests for 2024 events at Montpelier and for Marshall will be submitted at
project initiation.
Qualitative themes are encoded as binary indicators for Subtask 1.4 feature
importance analysis; the same organizational contact records produce the
building-level exclusion indicator for the Subtask 3.1 network analysis, coded
identically across both uses via the shared street address identifier.

\textbf{Subtask 1.4: Feature importance analysis of recovery determinants (lead: Chi, Napolitano).}
Building on the machine learning methodology validated in the Montpelier preliminary
work across 271 heritage structures~\cite{maalouf2025a,kallas2024}, Chi and Napolitano
will conduct feature importance analysis using random forest classifiers and
multi-layer perceptron models with two distinct outcome variables: collective action
participation status and NPS intervention typology classification.
The first analysis predicts collective action participation from network position,
institutional capacity, geographic location within the district, and socioeconomic
variables, directly addressing RQ1 and providing the explanatory layer for the
network and organizational analysis in RQ3.
The second analysis predicts NPS intervention typology classification from the same
feature set plus collective action participation status as a predictor, characterizing
the relative importance of participation versus building and owner characteristics
in explaining intervention outcomes.
This analysis is explicitly descriptive: collective action participation and
NPS-consistent outcomes are produced by the same organizational process and are
not causally independent; the analysis asks whether participation remains an
important predictor after controlling for building characteristics the owner cannot
change, which (if true) strengthens the equity argument for inclusive outreach.
Hierarchical correlation clustering will be implemented before importance analysis
to mitigate bias from correlated features~\cite{maalouf2025a}, and partial dependence
analysis will characterize the direction and magnitude of influential features on
each outcome.
The two sites present different temporal recovery depths: Montpelier's July 2023
flood and Marshall's September 2024 Hurricane Helene recovery are at structurally
different stages.
Rather than treating this as a confound to adjust away, separate site-specific
models are run at each site and feature importance rankings are compared directly.

The feature set combines binary indicators from NVivo thematic
coding~\cite{boyatzis1998transforming} (preservation priority, urgency, institutional
capacity, collective action exposure, financial constraint, expertise access, awareness
of district-scale intervention options), building and owner characteristics, and
Census data~\cite{uscensus_acs5year}.
Binary encoding avoids false precision in scaling subjective interview content.
Outcome variables are modeled separately to keep interpretation clean.
Cross-validation uses 5-fold at Montpelier (271 buildings) and leave-one-out
at Marshall where the smaller sample warrants maximum use of training data.
At Marshall, feature set size is capped at 15 variables prior to classifier
training, selected through recursive feature elimination to maintain a viable
feature-to-sample ratio and limit overfitting risk given the smaller district.
\textbf{Key outcomes and evaluation:}
RO1 produces four linked outputs that constitute the coupled dataset RO2 and RO3 require.
Subtask 1.1 delivers a building-level database for all structures at both sites,
each record carrying an NPS intervention typology classification at both building
and district scales and collective action participation status linked by street
address to all downstream data, complete by month 12.
Subtask 1.2 delivers calibrated HEC-RAS 2D baseline models at 1~m resolution for
both sites, with calibration validated against observed flood extents and
district-scale intervention states before RO2 begins.
Subtask 1.3 delivers 90--110 coded interview transcripts with codes converted to
binary indicators, collective action exposure linked per building, and awareness
of district-scale intervention options documented, continuing until data saturation
is confirmed.
Subtask 1.4 delivers site-specific feature importance models for both collective
action participation and NPS intervention typology, with convergent and divergent
rankings across sites characterizing which features are stable versus
context-dependent, contingent on inter-rater reliability achieving $\kappa \geq 0.70$.

% -----------------------------------------------------------------------
\begin{tcolorbox}[colback=rootwocol!60,colframe=black,boxrule=1pt,coltext=black,
  left=1mm,right=1mm,top=1mm,bottom=1mm]
\textbf{Research Objective 2 (RO2, leads: Busse \& Napolitano): Test how individual
building hardening, district-scale physical interventions, and mixed approaches
compare in their effects on district-scale flood vulnerability, and whether any
intervention approach produces threshold behavior in collective flood outcomes.}
\end{tcolorbox}
RO2 addresses RQ2 and constitutes the core engineering systems contribution of the
proposal: the first empirical comparison of individual building hardening,
district-scale physical interventions such as riparian vegetation management and
drainage improvements, and mixed approaches in terms of their collective flood
vulnerability outcomes in URM heritage districts.
Busse produces district-scale hydrodynamic outputs across empirically documented
and synthetic infrastructure configurations from RO1; Napolitano applies spatial
interdependence analysis and unsupervised feature learning to those outputs to
characterize how intervention approach and spatial configuration shape collective
flood outcomes, and to test whether any intervention approach produces threshold
behavior in those outcomes.

\textbf{Subtask 2.1: Infrastructure configuration development (lead: Napolitano, Busse).}
Napolitano and Busse will translate the RO1 typology classifications and collective
action participation patterns into spatially explicit district-scale infrastructure
configurations, varying along four dimensions: participation structure (coordinated,
individual, non-participating), geographic distribution of hardened buildings,
NPS intervention typology classification, and intervention scale (building-level
hardening, district-scale physical intervention, and mixed approaches).
Intervention typologies map to HEC-RAS 2D parameters: dry floodproofing as blocked
obstructions; wet floodproofing and no-action as flow-through with adjusted
Manning's $n$; terrain adaptations as DEM edits; fill-basement as reduced storage;
utility protection as hydraulic neutral; riparian vegetation as Manning's $n$
adjustments along the riverbank; drainage improvements and regrading as terrain
edits to the DEM.
Unknown buildings default to no-action with sensitivity flagging.
Observed configurations from RO1 establish baseline runs, including Montpelier's
empirically documented post-2023 riparian vegetation work as the sole observed
district-scale configuration; synthetic configurations vary hardening rate
(0--100\% in 10\% steps), spatial clustering (clustered/distributed/random),
participation structure (concentrated/distributed), and intervention scale
(building-level only, district-scale only, mixed).
Simulations will be executed on the Penn State Roar high-performance computing
cluster, where HEC-RAS 2D run times at 1~m resolution are estimated at
15--30~minutes per configuration, making a total of 500--1000 configurations
per site computationally feasible within the project timeline and providing
sufficient coverage of the configuration space for Subtask 2.3's comparative
analysis.

\textbf{Subtask 2.2: District-scale hydrodynamic modeling (lead: Busse).}
Busse will run district-scale hydrodynamic simulations across the configurations
developed in Subtask 2.1 at both sites, producing outputs of flood inundation
depth, flow velocity, and flood duration at the building level across the
heritage district.
Each simulation produces 1~m resolution raster outputs of flow velocity, inundation
depth, and water surface elevation across the full district.
These are aggregated to building-level scalar values using mean depth over each
building footprint, producing the paired configuration-outcome dataset that
Subtask 2.3 requires for the intervention comparison and threshold analysis.

All configurations are run at two return periods: the observed event magnitude
at each site (the July 2023 flood at Montpelier and Hurricane Helene at
Marshall, both directly calibrated) and the 100-year event (the standard design
flood and the most policy-relevant magnitude for pre-disaster investment
decisions~\cite{CFR44_59}).
The return period design balances coverage against computational feasibility.
A representative subset, comprising the empirically observed configurations plus
synthetic configurations at boundary participation rates (0\%, 50\%, and 100\% of
district buildings hardened) across all three intervention scales, are additionally
run at the 10-year and 50-year return periods to characterize how intervention
approach comparisons and any threshold behavior vary across flood magnitudes
without the full computational cost of all-configuration coverage at all return
periods.
The approximate total of 2,100 runs at 15--30 minutes per run on the Penn State
Institute for Computational and Data Sciences (ICDS) Roar cluster is feasible
within the project timeline, and PI Napolitano holds affiliate faculty status
with ICDS; she and PI Busse each have an existing annual computing allocation
that can be used for this.

Model uncertainty is characterized using two approaches.
Calibration residuals are reported against observed high water marks, post-event
survey data, and satellite-derived inundation extents for the July 2023 Montpelier
and Hurricane Helene Marshall events, quantifying model performance at the
building level before configurations are varied.
Manning's $n$ roughness coefficients for key land cover types are varied
$\pm$20\%, with sensitivity of mean building-level depth to roughness perturbation
reported across the configuration space; configurations where depth outputs are
sensitive to roughness assumptions are flagged in the Subtask 2.3 analysis.

\textbf{Subtask 2.3: Spatial interdependence and comparative intervention analysis (lead: Napolitano).}
Napolitano will apply exploratory data analysis and unsupervised feature learning
to the paired configuration-outcome dataset from Subtask 2.2 to characterize how
individual building hardening, district-scale physical interventions, and mixed
approaches compare in their effects on district-scale flood vulnerability, and to
test whether any intervention approach produces threshold behavior in collective
flood outcomes.
Results are compared across sites to assess whether intervention comparison findings
are consistent or context-dependent, generating propositions for future work rather
than generalizable claims from two cases.

The analysis proceeds in two stages that address distinct dimensions of the
configuration space.
First, intervention approach comparisons are drawn directly from the
configuration-outcome dataset: mean district-level flood depth and its distribution
across buildings are compared across building-level only, district-scale only, and
mixed configurations at equivalent participation rates, characterizing whether
district-scale intervention produces systematically different collective outcomes
than building-level hardening at the same level of community investment.
Second, PELT change-point detection~\cite{killick2012optimal} is applied within
each intervention class to the participation-rate dimension: mean district-level
flood depth as a function of aggregate hardening rate (0--100\% in 10\% steps),
which is a scalar ordered sequence suited to PELT input.
Discontinuities that survive permutation testing indicate threshold behavior;
smooth relationships indicate approximate linearity.
PCA-reduced configuration features are submitted to hierarchical agglomerative
clustering with Ward linkage and Euclidean distance~\cite{ward1963hierarchical}
to identify whether distinct outcome regimes exist across the full multidimensional
configuration space, including spatial pattern variation not captured by the
participation-rate curve alone.
Cluster count is selected by the largest gap in the dendrogram linkage distances,
avoiding the need to pre-specify the number of regimes.

Where threshold behavior is identified within any intervention class,
variance-based global sensitivity analysis (Sobol indices)~\cite{saltelli2010variance}
characterizes which buildings or district-scale features contribute most to variance
in collective flood outcomes across configurations.
Sobol indices are estimated using Gaussian process emulation fitted to the
500--1000 configuration-outcome pairs from Subtask 2.2, following established
metamodel-based sensitivity analysis procedures~\cite{saltelli2010variance,marrel2009calculations}.
The emulator evaluates sensitivity analytically from the fitted response surface,
requiring no additional model runs beyond those already planned.
Spatial autocorrelation of high-sensitivity locations using Moran's $I$ and LISA
tests whether they cluster spatially within the district or are distributed
independently of geography.
Findings are communicated to practitioner audiences in RO3 as intervention
comparison results and, where threshold behavior is identified, as maps of
high-sensitivity locations suitable for direct use in pre-disaster investment
planning.
If intervention approaches produce broadly similar outcomes, that finding is
reported as a null result and passed to RO3 as baseline physical exposure
(depth, velocity, duration from Subtask 1.2); the equity contribution of RO3
holds regardless of finding direction.

\textbf{Key outcomes and evaluation:}
RO2 produces two outputs that feed into RO3.
Subtask 2.1 delivers infrastructure configurations at both sites from RO1 typology
classifications across all three intervention scales plus synthetic configurations
across the full participation-rate range, archived by month 24.
Subtask 2.2 delivers building-level flood depth, velocity, and duration outputs
across all configurations at both return periods, with calibration validated against
observed extents and FEMA flood maps, and Manning's $n$ sensitivity analysis
complete before comparative analysis proceeds.
Subtask 2.3 delivers intervention comparison results across building-level,
district-scale, and mixed configurations, PELT-based threshold analysis with
permutation test p-values where threshold behavior is tested, and Sobol-based
sensitivity maps and LISA analysis where threshold behavior is identified, or
documented null findings passed to RO3 where they are not.
% -----------------------------------------------------------------------
\begin{tcolorbox}[colback=roothreecol!60,colframe=black,boxrule=1pt,coltext=black,
  left=1mm,right=1mm,top=1mm,bottom=1mm]
\textbf{Research Objective 3 (RO3, leads: Chi, Napolitano \& Busse): Characterize
the organizational and network structures that produce collective action exclusion
and explain why district-scale physical intervention was not pursued where technical
expertise identified it as an option, and translate findings into a field-tested
toolkit co-developed with APT DRI and the Department of the Interior.}
\end{tcolorbox}
RO3 addresses RQ3 by combining the collective action exclusion patterns and network
structures from RO1 with the intervention comparison outputs from RO2 to characterize
what organizational and network structures explain exclusion from collective action,
and whether those same structures explain why communities defaulted to fragmented
individual building hardening rather than district-scale intervention even where
technical expertise identified it as an option.
Toolkit outputs translate these findings into actionable guidance for pre-disaster
investment decisions, including intervention comparison results from RO2 and
exclusion pattern maps from RO1, each adapted for preservation specialists, property
owners, and community leaders.

\textbf{Subtask 3.1: Network and organizational analysis of collective action exclusion (lead: Chi, Napolitano).}
Chi and Napolitano will combine the building-level exclusion indicator established
in Subtask 1.3 with the intervention comparison outputs from RO2 to characterize
what organizational and network structures explain exclusion from collective action,
and whether those same structures explain why communities defaulted to fragmented
individual building hardening rather than district-scale intervention where technical
expertise identified it as an option.
Chi will lead the social network and institutional analysis that characterizes why
specific buildings and owners were excluded from collective action, framing exclusion
in terms of network position, institutional access, and geographic location within
the district~\cite{yin2021characterizing,yang2020untangle}; this analysis produces
the explanatory layer for the coordination failure pattern.
Napolitano will lead the join of Chi's exclusion characterization with the RO2
intervention comparison outputs, testing whether the organizational dynamics that
produce exclusion also predict which intervention approach a community pursues,
and characterizing the degree and structure of any correspondence.
Analysis will also identify which property attributes (geographic position within
the district, ownership type, social network connectivity) predict exclusion from
collective action, and assess whether the exclusion pattern and its relationship
to intervention approach is consistent across Montpelier and Marshall or differs
in ways that generate hypotheses for future work.
Findings will be reported as a set of testable propositions about the relationship
between collective action exclusion and intervention approach in heritage districts.

The network has two node types (property owners and organizational actors including
nonprofits (Montpelier Alive, Downtown Marshall Association), local government
(Montpelier Commission for Recovery), state historic preservation offices (VT and
NC SHPO), and federal agencies), with edges from organizational contact records
and interview data.
Three standard measures per owner node (degree centrality, betweenness centrality,
core-periphery structure) are complemented by institutional access indicators
(FEMA assistance, SHPO grants, prior heritage organization relationships) and
awareness of district-scale intervention options documented in Subtask 1.3.
These are linked to each building via street address and delivered as a
georeferenced table that Napolitano joins to the RO2 outputs using the shared
identifier.

Cross-case comparison tests whether the exclusion pattern (which buildings are
peripheral, which are core-adjacent) is consistent or different across two sites
that differ in network density, institutional density, and program
access~\cite{chi2012impacts}.
Findings are reported as propositions rather than causal claims: for example,
whether peripheral network position predicts exclusion more strongly in the
less institutionally dense site (Marshall) than in the more institutionally
supported one (Montpelier), and whether awareness of district-scale intervention
options follows the same network position pattern as collective action participation.
Consistent findings across both sites provide a stronger basis for future
comparative work than either site alone; inconsistent findings identify
context-dependence worth investigating at larger scale.

\textbf{Subtask 3.2: Stakeholder workshops and co-production (lead: Napolitano).}
Napolitano will organize in-person workshops in Montpelier and Marshall, followed
by a joint virtual cross-site session.
Each in-person workshop spans two days: day one presents RO1 and RO2 findings
with intervention comparison results and exclusion patterns, and day two uses
affinity diagramming to co-produce ranked pre-disaster collective action strategies
as primary toolkit inputs.
Participant lists are drawn from the RO1 exclusion data to ensure previously
excluded owners are present; outreach to this group uses direct mail and
door-to-door contact by local community liaisons at Montpelier Alive, the
Montpelier Commission for Recovery, and the Downtown Marshall Association already
engaged through PI Napolitano's APT DRI relationships at both sites, rather than
relying on the organizational networks that originally failed to reach them.
Wellock will review outputs for federal alignment before final toolkit production.

\textbf{Subtask 3.3: Toolkit development and testing (lead: Napolitano).}
Drawing on comparative analysis findings from Subtask 3.1 and the co-produced
strategies from Subtask 3.2, Napolitano, who directs the APT DRI, will develop
the Flood-Resilient Heritage Preservation Toolkit through APT DRI.
The toolkit serves three audiences: preservation specialists with NPS-grounded
flowcharts and pre-disaster assessment templates calibrated to the four intervention
windows~\cite{maalouf2025a,maalouf2025b}; property owners with tip sheets modeled
on FLASH~\cite{FLASH-SafeRoomChecklist} and FEMA resources~\cite{FEMA-ProtectYourProperty};
and community leaders with intervention comparison results, exclusion pattern maps,
and guidance on organizational structures that enable district-scale collective action.
Communication components draw on the Natural Hazards Center's Principles of Risk
Communication~\cite{campbell2021}, adapting their guidance on communicating through
trusted messengers, providing clear and actionable information, and tailoring
messages for target audiences to the specific constraints of URM heritage district
owners navigating the preservation-urgency paradox.
Testing at the annual APT conference tabletop exercise measures whether practitioners
make better pre-disaster decisions using the toolkit, scored against a benchmark
established by an expert panel of five practitioners: two preservation engineers
from APT DRI, at least one state historic preservation officer with flood recovery
experience (VT or NC), and Wellock as the federal heritage recovery representative,
all of whom PI Napolitano has worked with before.
Wellock's Department of the Interior connections extend reach to heritage-at-risk
communities nationwide.
All components are published openly through the APT DRI platform and archived on
NHERI DesignSafe with a CC BY 4.0 license.
Tabletop exercise success is measured through a pre/post questionnaire design:
participants make a building prioritization decision before and after using the
toolkit, with decision quality scored against an expert panel benchmark from the
RO2 and RO3 findings.

Dissemination beyond the two case study sites operates through two channels.
A three-part webinar series delivered through APT reaches preservation
professionals, state historic preservation officers, and local emergency managers
across the APT network, with sessions covering the research findings, the toolkit
components, and the workshop co-production model for communities wanting to adapt
the toolkit to their own contexts.
Wellock's connections to heritage-at-risk communities through the Department of
the Interior provide a second dissemination pathway, reaching communities already
identified through federal programs; the DOI's established coordination with HENTF
extends this reach to cultural stewards embedded in local emergency management
infrastructure nationwide, giving the toolkit direct access to the federal heritage
recovery infrastructure it was designed to inform.
Risk communication strategy draws on resources from Katelyn Nolan (USACE) and the
Natural Hazards Center to ensure toolkit messaging reaches emergency management
and floodplain management audiences beyond the preservation community.

\textbf{Key outcomes and evaluation:}
RO3 produces three linked outputs.
Subtask 3.1 delivers testable propositions about the organizational and network
structures that produce collective action exclusion and explain the failure to
pursue district-scale intervention at each site, with the georeferenced join table
archived on DesignSafe and reviewed before stakeholder presentation by month 27.
Subtask 3.2 delivers ranked pre-disaster collective action strategies co-produced
in workshops at both sites and a joint virtual session, with participant lists
drawn from the RO1 exclusion data to ensure previously excluded owners are present.
Subtask 3.3 delivers the Flood-Resilient Heritage Preservation Toolkit published
through APT DRI and archived on DesignSafe under CC BY 4.0 by month 36, with
field testing measured through pre/post questionnaire data from the APT conference
tabletop exercise.
\vspace{-0.4cm}
% -----------------------------------------------------------------------
\customsection{Broader Impacts}
% -----------------------------------------------------------------------
\vspace{-0.2cm}
This research provides the first empirical comparison of individual building
hardening, district-scale physical intervention, and mixed approaches in terms
of their collective flood vulnerability outcomes in URM heritage districts, and
the first characterization of the organizational and network structures that
explain why communities default to fragmented individual hardening rather than
coordinated district-scale intervention.
Either finding gives communities a quantitative foundation for pre-disaster
investment that currently has none, delivered through the Flood-Resilient Heritage
Preservation Toolkit with Wellock's Department of the Interior connections extending
reach to heritage-at-risk communities nationwide.
The social-network-plus-hydrodynamic-modeling methodology is transferable to other
hazards where individual and collective intervention decisions aggregate to
infrastructure-scale outcomes, and is replicable regardless of findings at
Montpelier and Marshall.
Where findings show district-scale intervention produces better collective flood
outcomes, the organizational analysis identifies the social levers communities
would need to shift to pursue it; where findings show intervention approaches
produce comparable outcomes, that establishes a baseline and suggests outreach
need not be redirected toward district-scale coordination.
A future research direction this work opens is characterizing what institutional
conditions and community capacities would need to be in place for communities to
pursue district-scale intervention at scale, a question the current research
motivates but does not fully answer.

This research advances the ISP program's interest in infrastructure system
function under extreme events and the equity dimensions of infrastructure
performance, contributing to NEHRP and NWIRP goals: URM buildings were identified as critical infrastructure vulnerabilities in the NEHRP seismic
context, and the hurricane- and storm-driven flooding studied here connects
directly to NWIRP's interest in windstorm impacts on communities.
It extends the NIST SP 1310 socio-technical systems framing to flood hazards
and heritage infrastructure, responding directly to NIST SP 1310's call for
analogous frameworks for hurricanes and floods~\cite{NIST2024SP1310} and the
FEMA/BSSC Functional Recovery Planning Committee's parallel call for retrofit
strategies for existing construction~\cite{NIBSFunctionalRecovery2023}.
The project leverages NSF's investment in NHERI by archiving all field data,
hydrodynamic simulation outputs, and georeferenced datasets in the NHERI
DesignSafe Data Depot under open licenses, building on PI Napolitano's prior
NHERI RAPID Graduate Scholar datasets already archived there.
Findings will be published in the International Journal of Disaster Risk Reduction,
Journal of Cultural Heritage, and Natural Hazards and Earth System Sciences, and
presented at Engineering Mechanics Institute, APT, and Keeping History Above Water
conferences.

Research outcomes integrate into education through PI Napolitano's Penn State
courses (200-level programming for 120 architectural engineers annually; 400-level
structural engineering for 40 students) and PI Chi's Indiana University courses
(300-level population geography; 400- and 500-level spatial analysis).
Post-doctoral scholars develop domain-specific skills: Napolitano's lab in spatial
data science and community engagement; Busse's lab in district-scale hydrodynamic
modeling; Chi's lab in mixed-methods social science and network analysis.
Undergraduates in Built Environment Seminal Undergraduate Research Experience
(BESURE), WISER (Women in Science and Engineering Research), and MURE (Minority
Undergraduate Research Experience) contribute to building typology classification,
geospatial data preparation, and HEC-RAS 2D output post-processing at no cost to
the project; PI Napolitano, serving as undergraduate research coordinator for her
department and BESURE, has direct access to department-sponsored undergraduate
researchers and existing program infrastructure for recruiting, supervising, and
tracking their contributions.
Learning outcomes are assessed through pre/post surveys, faculty rubrics, and
long-term STEM career tracking.
All code will be released on GitHub under an MIT license and datasets deposited
in the NHERI DesignSafe Data Depot~\cite{maalouf2024data,maalouf2026data}.
Toolkit materials will be freely accessible through the APT DRI platform and
distributed through local government channels, community organizations, and
social media.
\vspace{-0.4cm}
% -----------------------------------------------------------------------
\section{Project timeline and evaluation}The project runs from January 2027 through December 2029 (36 months, Figure~\ref{fig:gantt}).
% -----------------------------------------------------------------------
\begin{figure}[H]
\centering
\resizebox{0.8\textwidth}{!}{%
\begin{ganttchart}[
  hgrid,
  vgrid=false,
  x unit=2.6mm,
  y unit title=5mm,
  y unit chart=4.8mm,
  title/.style={fill=lightgray, draw=black, font=\scriptsize\bfseries},
  title label font=\scriptsize\bfseries,
  bar label font=\scriptsize,
  group label font=\scriptsize\bfseries,
  bar height=0.55,
  group height=0.3,
  bar top shift=0.225,
  bar label node/.append style={align=right, text width=4cm},
  group label node/.append style={align=right, text width=4cm},
  milestone label font=\tiny,
  milestone label node/.append style={text width=0pt, text=white},
]{1}{36}
%--- title rows
\gantttitle{Year 1 (2027)}{12}\gantttitle{Year 2 (2028)}{12}\gantttitle{Year 3 (2029)}{12} \\
\gantttitle{Q1}{3}\gantttitle{Q2}{3}\gantttitle{Q3}{3}\gantttitle{Q4}{3}%
\gantttitle{Q1}{3}\gantttitle{Q2}{3}\gantttitle{Q3}{3}\gantttitle{Q4}{3}%
\gantttitle{Q1}{3}\gantttitle{Q2}{3}\gantttitle{Q3}{3}\gantttitle{Q4}{3} \\
%--- RO1
\ganttgroup[
  group/.append style={fill=roonecol, draw=roonecol!70!black},
  group label node/.append style={text=roonecol!70!black}
]{RO1}{1}{18}
  \ganttmilestone[milestone/.append style={fill=roonecol!60!black, draw=roonecol!40!black, shape=rectangle, rounded corners=1pt, inner sep=6pt, label={[font=\tiny\bfseries,text=white,inner sep=0pt]center:M1}}]{}{12}
  \ganttmilestone[milestone/.append style={fill=roonecol!60!black, draw=roonecol!40!black, shape=rectangle, rounded corners=1pt, inner sep=6pt, label={[font=\tiny\bfseries,text=white,inner sep=0pt]center:M2}}]{}{18} \\
\ganttbar[bar/.append style={fill=roonecol!55, draw=roonecol!80}]{1.1 Typology \& CA status}{1}{12} \\
\ganttbar[bar/.append style={fill=roonecol!55, draw=roonecol!80}]{1.2 Site characterization}{1}{12} \\
\ganttbar[bar/.append style={fill=roonecol!55, draw=roonecol!80}]{1.3 Interviews}{3}{18} \\
\ganttbar[bar/.append style={fill=roonecol!55, draw=roonecol!80}]{1.4 Feature importance}{12}{18} \\
%--- RO2
\ganttgroup[
  group/.append style={fill=rootwocol, draw=rootwocol!70!black},
  group label node/.append style={text=rootwocol!70!black}
]{RO2}{12}{27}
  \ganttmilestone[milestone/.append style={fill=rootwocol!60!black, draw=rootwocol!40!black, shape=rectangle, rounded corners=1pt, inner sep=6pt, label={[font=\tiny\bfseries,text=white,inner sep=0pt]center:M3}}]{}{27} \\
\ganttbar[bar/.append style={fill=rootwocol!55, draw=rootwocol!80}]{2.1 Configurations}{12}{18} \\
\ganttbar[bar/.append style={fill=rootwocol!55, draw=rootwocol!80}]{2.2 Hydrodynamic modeling}{12}{24} \\
\ganttbar[bar/.append style={fill=rootwocol!55, draw=rootwocol!80}]{2.3 Comparative analysis}{18}{27} \\
%--- RO3
\ganttgroup[
  group/.append style={fill=roothreecol, draw=roothreecol!70!black},
  group label node/.append style={text=roothreecol!70!black}
]{RO3}{18}{36}
  \ganttmilestone[milestone/.append style={fill=roothreecol!60!black, draw=roothreecol!40!black, shape=rectangle, rounded corners=1pt, inner sep=6pt, label={[font=\tiny\bfseries,text=white,inner sep=0pt]center:M4}}]{}{30}
  \ganttmilestone[milestone/.append style={fill=roothreecol!60!black, draw=roothreecol!40!black, shape=rectangle, rounded corners=1pt, inner sep=6pt, label={[font=\tiny\bfseries,text=white,inner sep=0pt]center:M5}}]{}{35} \\
\ganttbar[bar/.append style={fill=roothreecol!55, draw=roothreecol!80}]{3.1 Network analysis}{18}{27} \\
\ganttbar[bar/.append style={fill=roothreecol!55, draw=roothreecol!80}]{3.2 Workshops}{24}{27} \\
\ganttbar[bar/.append style={fill=roothreecol!55, draw=roothreecol!80}]{3.3 Toolkit development}{27}{33} 
\end{ganttchart}%
}
\vspace{-0.3cm}
\caption{Project timeline. Rectangles (M1--M5) indicate key performance indicators.}
\label{fig:gantt}
\end{figure}
\vspace{-0.3cm}

\newpage 
RO1 runs concurrently in years 1--2: Subtasks 1.1--1.2 begin month 1; Subtasks
1.3--1.4 span months 3--18.
RO2 begins month 12 once RO1 typology and site characterization are sufficiently
developed, with modeling through month 24 and comparative analysis complete by month 27.
RO3 runs months 18--36: network and organizational analysis months 18--27, workshops
months 24--27, toolkit development months 27--33, APT testing months 33--36
(Fig.~\ref{fig:gantt}).
Progress will be evaluated against six key performance indicators corresponding to
the three research objectives:
(1) Completion of NPS intervention typology classification at both building and
district scales and collective action status recording for 271 structures in
Montpelier and over 70 structures in Marshall (RO1, month 12);
(2) Completion of social science data collection with target interview sample sizes,
data saturation confirmed, and awareness of district-scale intervention options
documented (RO1, month 18);
(3) Delivery of calibrated district-scale hydrodynamic models for both sites (RO2,
month 18);
(4) Completion of hydrodynamic simulations across documented configurations and
intervention comparison analysis, with threshold behavior reported where present
(RO2, month 27);
(5) Completion of stakeholder workshops in Montpelier and Marshall with
co-produced toolkit draft (RO3, month 30);
(6) Field-tested toolkit released on APT DRI platform following APT conference
tabletop evaluation (RO3, month 36).
Progress will be monitored through weekly project meetings and tracked in PSU
Microsoft Planner and Teams.
\vspace{-0.3cm}
% -----------------------------------------------------------------------
\section{Results from prior NSF Support}
% -----------------------------------------------------------------------

\textbf{PI Napolitano:}
CMMI-2222849, \$131,317, April 2022 -- February 2025.
RAPID: Examining the Performance of Historic Masonry Structures after the Midwest Tornadoes.
Intellectual Merit: This project addressed building code gaps by analyzing historic masonry structures' performance under tornadic loading and assessing retrofitting efficacy, producing the first systematic dataset on tornado damage to unreinforced masonry heritage structures.
Broader Impacts: The project enhanced understanding of historic masonry resilience under wind hazards, produced open datasets archived on DesignSafe, and integrated community-centered structural data into engineering education.
Publications: \cite{kaushal2023understanding, kaushal2023structural, kaushal2023b, sskmgeometric, sskmdatacollection, sskmstructuralanalysis}.
Undergraduate research presentations: \cite{patelski2022research, whitesides2023research}.
Publicly shared datasets: \cite{steerquad, MapillaryDatasetSteer360, kaushal2023processed, kaushal2023raw}.
Renewal: N/A.

\noindent\textbf{Co-PI Busse:}
% maggie: add prior NSF support in the same format above.
% If no prior NSF support, note relevant prior funding and publications
% that establish hydrodynamic modeling expertise.
is senior personnel on NSF CBET-2400672, \$1,000,000, August 2024 - July 2029:
RAISE: CET: Co-development of engineering energy literacy and engineering identities through campus decarbonization research.
Intellectual Merit:
This work integrates six component-specific research directions to address fundamental questions associated with net-zero fuels in the context of Combined Heat and Power decarbonization.
Intellectual Merit of the research lies in the integrated real-world system evaluation approach as well as the evaluation of student energy literacy through their research experience. 
Broader Impacts:
Advancement of decarbonization technologies and the development of an energy-literate future engineering workforce.  
Publications: None.
Undergraduate research presentations: \cite{little2026}
Renewal: N/A

\noindent\textbf{Co-PI Chi:}
ICER-1927827, \$3,018,000, January 2020 -- June 2025.
NNA Track 1: Pursuing Opportunities for Long-term Arctic Resilience for Infrastructure and Society (POLARIS).
Intellectual Merit: This transdisciplinary project developed theoretical frameworks and models addressing environmental disruption to communities and infrastructure, food security, and migration in Arctic Indigenous communities, producing findings applicable to other resource-constrained communities facing existential environmental threats.
Broader Impacts: Findings inform government, researcher, practitioner, and Indigenous group decision-making for adaptation to rapid Arctic change; the project supported 6 Ph.D.\ students, 3 postdocs, and 10+ undergraduates.
Publications: 16 peer-reviewed articles \cite{yin2022tale, korkut2022addressing, Yin2022, Holen2022, Mucioki2024, Chi2024, Christian2024, Smith2024, Abdar2020, Holen2023, Mucioki2024b, Zhou2024, Mucioki2024c, Hamilton2024, Smith2023, Powell2023}, 6 op-ed pieces \cite{ArcticPOLARIS, Bogardus2024, Baldwin2023, AumanBauer2019, POLARIS, AumanBauer2020}, 5 videos \cite{WoodRiverBridge2022, StakesForStakeholders2022, POLARISTeam2021, ChevakSchool2024, ChevakAlaska2022}, and 16 conference presentations \cite{Smith2025, Maio2024, Yin2025, Maio2024b, Miller2024, Miller2021, Mucioki2024, Yin2024, Tickamyer2024, Tickamyer2025, Hamilton2024, Maio2024c, Chi2024, Mucioki2024b, Chi2025, Maio2024d}.
Renewal: N/A.
\newpage
\bibliography{biblio}

\end{document}
